\documentclass{article}
\usepackage{amsmath,amssymb}
\usepackage[vocab={designs}]{leanify}
\newtheorem{lemma}{Lemma}
\begin{document}

\section{Notation}

We write $r(v,k,\lambda)$ for the replication number of a design with $v$ points, block size $k$ and index $\lambda$.
We write $b(v,k,\lambda)$ for the number of blocks of a design with $v$ points, block size $k$ and index $\lambda$.
We write $T(v)$ for the number of triples of a Steiner triple system of order $v$.
We write $P(n)$ for the number of points of a projective plane of order $n$.
We write $L(n)$ for the number of lines of a projective plane of order $n$.
We write $A(n)$ for the number of points of an affine plane of order $n$.
We write $B(n)$ for the number of lines of an affine plane of order $n$.

\section{The Fano plane}

\begin{lemma}\label{fano}
$r(7,3,1)=3$, $b(7,3,1)=7$, $T(7)=7$, $P(2)=7$ and $L(2)=7$.
\end{lemma}

\begin{lemma}\label{affine}
$A(3)=9$, $B(3)=12$ and $T(9)=12$.
\end{lemma}

\begin{lemma}\label{admissible}
$7$ is admissible for Steiner triple systems and $9$ is admissible for Steiner triple systems.
\end{lemma}

\section{Counting}

\begin{lemma}\label{count}
Let $v\ge7$ be an integer that is admissible for Steiner triple systems. Then $T(v)=\frac{v(v-1)}{6}$ and $3T(v)=v\,r(v,3,1)$.
\end{lemma}

\begin{lemma}\label{planes}
$P(3)=13$, $L(3)=13$ and $B(4)=20$.
\end{lemma}

\end{document}
